% sections/B_preprocessing.tex: Lane p1.
% Owner: the writer named in the report plan. Numbers only through \R / \Rci; figures through \genfig,
% tables through \gentab (report/README.md).
\section{Preprocessing and data-treatment decisions}\label{app:prep}

This appendix justifies every signal-processing and data-treatment choice of the pipeline (lane p1).
Each choice was tested by changing one setting at a time and re-running the full leave-one-person-out
(LOSO) evaluation with a fixed panel of three models.
Every default survives this test. The appendix also covers the calibration of the panel's \pright{} and
the combination of several attempts.
Abbreviations: confidence interval (CI), logistic regression (LR), gradient-boosted trees (GBM),
MiniRocket with a ridge classifier (MR), area under the ROC curve (AUC), expected calibration error (ECE),
Benjamini--Hochberg (BH). A difference in accuracy is given in percentage points (points).

% ------------------------------------------------------------------------------------------------
\subsection{How each choice was judged}\label{app:prep:method}

\paragraph{The panel.} Three models of different kinds were fixed before any option was compared, each
with one configuration (Table~\ref{tab:p1_panel}).
\emph{LR}: an L2 logistic regression ($C = \R{p1/lr_c}$) on the means of \R{p1/n_bins} time bins of the
\R{p1/lr_n_channels} chiral gravity-frame channels (the yaw rate $\vomega\cdot\ghat$ and its cumulative
integral), trained with mirror augmentation.
With mirror augmentation the fitted LR is antisymmetric, $p(m(x)) = 1 - p(x)$, so mirror-even inputs get no
weight (Section~\ref{sec:physics}).
\emph{GBM}: LightGBM with fixed parameters (\R{p1/gbm_n_estimators} trees, learning rate
\R{p1/gbm_learning_rate}, \R{p1/gbm_num_leaves} leaves, row and column subsampling \R{p1/gbm_subsample} and
\R{p1/gbm_colsample}) on \R{p1/gbm_features_per_channel} summary features of each of the
\R{p1/gbm_n_channels} channels, with mirror augmentation.
\emph{MR}: MiniRocket~\cite{dempster2021minirocket} (\R{p1/mr_features} features requested) on the
\R{p1/mr_n_channels} gravity-frame channels, with a ridge classifier whose penalty and Platt
scaling~\cite{platt1999probabilistic} were fitted by an inner \R{p1/mr_inner_splits}-fold split over the
training persons.
The panel spans a linear model on two physically motivated channels, a tree ensemble on engineered features
and a random-kernel time-series model. A treatment that helps or hurts all three is a property of the data,
not of one model.

\paragraph{Options.} An \emph{option} is an alternative to one default. Options on the recorded stream
(cut-off, filter order, causal filter, gap bridging, interpolation, sampling rate, rounding, timestamp
jitter, gravity estimate, window, bound errors) were built by re-standardising the raw records with one
setting changed. This produced \R{p1/n_variants} variant caches, including the default one. With every
setting at its default, the rebuilt arrays equal the analysis cache exactly (largest difference
\R{p1/variant_fidelity_max_diff} over \R{p1/variant_fidelity_n_trials} trials). Options on the training set
(channel sets, normalisation, winsorising, augmentation, training-row filters, platform) were applied inside
each training fold. In all, the lane ran \R{p1/n_experiments} LOSO experiments on the \R{p1/n_trials}
primary trials of \R{p1/n_persons} persons.

\paragraph{Paired comparison.} Each option was compared with the same model under the current pipeline,
on the same trials and folds. The comparison gives the pooled accuracy difference with its CI
(\R{p1/n_boot} person resamples), the mean per-person log-loss difference with its CI, and an exact
sign-flip test over persons. The comparisons form three families, each adjusted by BH and by Holm: the
options (\R{p1/fam_options_size} option-by-model comparisons), the robustness costs
(\R{p1/fam_robustness_size} test-time segmentation errors) and the pairs (\R{p1/fam_pairs_size} remedies
and size-matched controls).

\paragraph{Verdicts.} \emph{Worse} or \emph{better}: the paired CI excludes 0 and the BH $q$ is below
\R{p1/alpha}. \emph{Nominally worse} or \emph{nominally better}: only the CI excludes 0; this is suggestive,
not significant. \emph{Identical}: every prediction is within \R{p1/identical_tol} of the default (solver
precision). \emph{No difference}: the CI includes 0. This is absence of evidence, not equivalence. The
lower end of its CI is the largest loss the data allow, and the upper end is the largest gain forgone.

\paragraph{Decision rule.} The rule was fixed in advance: change a default only if one option is better for
at least two panel models and worse for none. An option run on one model cannot meet the rule; this
applies to one option of the frame choice and one of the augmentation choice.

\paragraph{Other roles.} A \emph{robustness} run perturbs the test input of an unchanged model; it sets a
requirement for production and is never an option. A \emph{remedy} is judged only against the problem it
addresses. A \emph{scoring} run scores the trials an exclusion rule removes with the base model. A
\emph{control} removes as many training trials or persons at random as an option does. The leak check
ran on every fold for the base and nested-calibrated runs, and passed; for the other variants it ran on
\R{p1/leak_sample_folds} seeded folds.

\gentab{p1/p1_panel}

The panel's natural-use result is right-hand recall on the free-choice trials (all right-hand so far; a
constant right answer would score 100\,\%): LR \R{p1/lr_nat_answered}, GBM \R{p1/gbm_nat_answered} and MR
\R{p1/mr_nat_answered} of \R{p1/gbm_nat_n} trials (Table~\ref{tab:p1_panel}). No panel model returned a
\pright{} of exactly 0.5 on these trials (coverage \R{p1/lr_nat_coverage}, \R{p1/gbm_nat_coverage} and
\R{p1/mr_nat_coverage}), so the forced-choice accuracy equals the answered accuracy.
GBM has the highest base accuracy, \Rci{p1/gbm_accuracy}. Calling it the best panel model is a choice made
on the outer folds among \R{p1/best_n_configs_tried} models, so it is exploratory.

% ------------------------------------------------------------------------------------------------
\subsection{Overview: every default is kept}\label{app:prep:overview}

Figure~\ref{fig:P1_01_forest} shows every option against the current pipeline for each panel model.
Table~\ref{tab:p1_decisions} summarises the \R{p1/n_choices} choices and Table~\ref{tab:p1_decisions_by_model}
gives the same choices per model. The rule keeps the default in \R{p1/n_choices_kept} of \R{p1/n_choices}
choices and changes none (\R{p1/n_choices_changed}). Unless stated otherwise, $q$ below is the BH-adjusted
value within the options family.

No option is better than the current pipeline, even nominally (\R{p1/n_options_better} better,
\R{p1/n_options_nominally_better} nominally better). Of the \R{p1/fam_options_size} option comparisons,
\R{p1/n_options_ci_excludes_0} have a CI that excludes 0; about \R{p1/expected_false_ci} would do so by
chance at the 95\,\% level. BH keeps \R{p1/n_options_worse}, both worse, and Holm keeps
\R{p1/n_options_holm}:
\begin{itemize}[nosep]
\item device-frame channels instead of the chiral channels in LR: \Rci{p1/d_lr_frame_device}~points,
  BH $q$~=~\R{p1/q_lr_frame_device}, Holm $p$~=~\R{p1/holm_lr_frame_device};
\item gravity-frame channels only in GBM: \Rci{p1/d_gbm_frame_gravity}~points,
  BH $q$~=~\R{p1/q_gbm_frame_gravity}, Holm $p$~=~\R{p1/holm_gbm_frame_gravity}.
\end{itemize}
The other \R{p1/n_options_nominally_worse} CIs that exclude 0 are nominally worse and do not survive either
adjustment. Of the remaining comparisons, \R{p1/n_options_identical} are identical to the default and
\R{p1/n_options_no_difference} show no difference.

The per-person log loss points the same way. Its paired CI lies above 0 for the device-frame LR
(\Rcib{p1/dll_lr_frame_device}), the gravity-frame GBM (\Rcib{p1/dll_gbm_frame_gravity}) and the GBM
without mirror augmentation (\Rcib{p1/dll_gbm_no_mirror}), among others.
A few options lower the log loss nominally without changing accuracy: the fixed 3-s window for MR
(\Rcib{p1/dll_mr_win_fixed}), the approach plus hold for GBM (\Rcib{p1/dll_gbm_win_hold}) and time warping
for GBM (\Rcib{p1/dll_gbm_time_warp}). The rule decides on accuracy, so these remain leads for future work,
not changes.
Figure~\ref{fig:P1_01_forest} shows the accuracy differences of all \R{p1/fam_options_size}
option-by-model comparisons. The full numbers (accuracy, accuracy difference with its CI, BH $q$, log-loss
difference with its CI and verdict) are in \code{tables/p1/p1_options_a} and \code{p1_options_b} of the
analysis outputs; at more than a page each, they are too long to reproduce here.

\begin{figure}[p]
\centering
\genfig{p1/P1_01_forest}
\caption{Every data-treatment option against the current pipeline, per panel model: paired difference in
LOSO accuracy (points) with its 95\,\% CI (\R{p1/n_persons} persons; \R{p1/fam_options_size} comparisons in
one multiplicity family). Filled: CI excludes 0 and BH $q < \R{p1/alpha}$; light fill: CI alone excludes 0;
open circle: no difference; open square: identical predictions. The exclusion-rule rows add the excluded
trials to training and score the primary trials only. No option beats the default for two models, so every
default is kept.}
\label{fig:P1_01_forest}\label{fig:app:prep-forest}
\end{figure}

\gentab{p1/p1_decisions}
\gentab{p1/p1_decisions_by_model}

% ------------------------------------------------------------------------------------------------
\subsection{Low-pass filter: cut-off, order and phase}\label{app:prep:filter}

\paragraph{Options.} The default is a zero-phase Butterworth low-pass at \R{p1/lowpass_hz}~Hz, order
\R{p1/lowpass_order}, applied after resampling to \R{p1/fs_hz}~Hz. The options were no low-pass, cut-offs
of 5, 10 and 20~Hz, orders 2 and 8, and a causal (one-pass) filter of the same design.

\paragraph{Evidence.} The movement lives below the cut-off: \R{describe/gyro_power_below_lowpass} of the
gyroscope power and \R{describe/acc_power_below_lowpass} of the accelerometer power of the approach lie
below 15~Hz, on both platforms (Figure~\ref{fig:S2_07_spectrum}). No filter option changes accuracy beyond
noise (Figure~\ref{fig:P1_02_cutoff}). Only the 5-Hz cut-off in LR is nominally worse,
\Rci{p1/d_lr_lp5}~points (BH $q$~=~\R{p1/q_lr_lp5}). The largest loss compatible with the data is
\Rb{p1/d_mr_lp_none_lo}~points (no low-pass, MR), and the largest gain forgone is
\Rb{p1/d_gbm_lp5_hi}~points (5~Hz, GBM).

\paragraph{Phase.} A causal filter delays the signal. At the current design it delays a 1-Hz movement by
\R{p1/causal_delay_ms}~ms; a 10-Hz cut-off would cost \R{p1/delay_10hz_o4_ms}~ms, a 20-Hz cut-off
\R{p1/delay_20hz_o4_ms}~ms, and order 8 at 15~Hz \R{p1/delay_15hz_o8_ms}~ms
(Figure~\ref{fig:P1_10_filter_delay}). Its accuracy cost is within noise: LR \Rci{p1/d_lr_lp_causal} and GBM
\Rci{p1/d_gbm_lp_causal}~points.

\paragraph{Decision.} Keep the zero-phase 15-Hz filter of order 4. The predictor processes a recorded
buffer, so a zero-phase filter is also available in production. If a later online stage needs a causal
filter, the data bound its accuracy cost at \Rb{p1/d_gbm_lp_causal_lo}~points.

\begin{figure}[htbp]
\centering
\begin{subfigure}[t]{0.49\linewidth}
\centering
\genfig{core/S2_07_spectrum}
\caption{Cumulative power spectrum of the approach (median over \R{describe/primary_trials} primary
trials); the vertical line is the 15-Hz cut-off.}
\label{fig:S2_07_spectrum}
\end{subfigure}\hfill
\begin{subfigure}[t]{0.49\linewidth}
\centering
\genfig{p1/P1_10_filter_delay}
\caption{Group delay of causal Butterworth low-pass filters (order 4) at \R{p1/fs_hz}~Hz; the zero-phase
filter of the pipeline has none.}
\label{fig:P1_10_filter_delay}
\end{subfigure}
\caption{Why the cut-off is 15~Hz and why the filter is zero-phase. (a)~Almost all movement power lies below
the cut-off on both platforms. (b)~A causal filter at the current design would delay a 1-Hz movement by
\R{p1/causal_delay_ms}~ms.}
\label{fig:app:prep-filter}
\end{figure}

\begin{figure}[htbp]
\centering
\genfig{p1/P1_02_cutoff}
\caption{Low-pass filter options. Left: cut-off of the zero-phase Butterworth filter (order 4); right: its
order at 15~Hz, and the causal filter of the same design. Paired accuracy difference against the current
pipeline (star), per panel model (legend: current accuracy), with 95\,\% CIs; markers as in
Figure~\ref{fig:P1_01_forest}. The default is kept.}
\label{fig:P1_02_cutoff}\label{fig:app:prep-cutoff}
\end{figure}

% ------------------------------------------------------------------------------------------------
\subsection{Gaps, sampling rate, timestamps and gravity}\label{app:prep:sampling}

\paragraph{Gaps and imputation.} The default bridges gaps of up to \R{p1/max_gap_ms}~ms by linear
interpolation. Trials with a longer gap near the movement are excluded by rule E4d
(Section~\ref{app:prep:exclusions}). Bridging gaps of up to 250~ms changes no primary trial
(\R{p1/vdiff_gap250_n_differing} trials differ), so its predictions are identical. Bridging only up to
50~ms changes \R{p1/vdiff_gap50_n_differing} trials and shows no difference. Cubic-spline and
sample-and-hold resampling show no difference either; the largest loss compatible with the data is
\Rb{p1/d_gbm_interp_hold_lo}~points (sample-and-hold, GBM). Decision: keep linear bridging up to
\R{p1/max_gap_ms}~ms.

\paragraph{Sampling rate and timestamps.} The predictor must tolerate slower browsers and imprecise
timestamps, because phones differ in rate and timing~\cite{stisen2015heterogeneity}. Keeping every second
or third raw sample (30 and 20~Hz), rounding the iOS values to 0.1, or adding Gaussian noise of SD 2, 5 or
10~ms to the sample times shows no difference for any model (Figure~\ref{fig:P1_04_sampling}). The largest
drop is LR at 20~Hz, \Rci{p1/d_lr_fs20}~points. The largest loss compatible with the data is
\Rb{p1/d_mr_fs20_lo}~points (20~Hz, MR). Decision: keep the native rate and precision, resampled to
\R{p1/fs_hz}~Hz. This is a bound, not a proof of robustness.

\paragraph{Gravity estimate.} The chiral channels are defined relative to gravity, so the gravity estimate
matters in principle. The default propagates gravity with the gyroscope. A 1-Hz low-pass of the
accelerometer and the browser's own fusion~\cite{w3c2025deviceorientation} show no difference
(Table~\ref{tab:p1_decisions}). The largest loss compatible with the data is
\Rb{p1/d_lr_grav_lowpass_lo}~points (1-Hz low-pass, LR); the largest gain forgone is
\Rb{p1/d_gbm_grav_browser_hi}~points (browser fusion, GBM). Decision: keep the gyroscope-propagated
estimate, because the browser fusion differs between platforms and is outside our control.

\begin{figure}[htbp]
\centering
\genfig{p1/P1_04_sampling}
\caption{Robustness to the recording. Paired accuracy difference against the current pipeline (star) when
every second or third raw sample is kept (30 and 20~Hz; the low-pass is capped at 0.4 times the rate), and
when Gaussian noise is added to the raw sample times. 95\,\% CIs; markers as in
Figure~\ref{fig:P1_01_forest}. A CI that includes 0 bounds the cost by its lower end; it does not show
equivalence.}
\label{fig:P1_04_sampling}\label{fig:app:prep-sampling}
\end{figure}

% ------------------------------------------------------------------------------------------------
\subsection{Representation: frame and window}\label{app:prep:frame}

\paragraph{Frame (channel set).} Each panel model has its own default: LR uses the
\R{p1/lr_n_channels} chiral gravity-frame channels, MR the \R{p1/mr_n_channels} gravity-frame channels and
GBM all \R{p1/gbm_n_channels} channels. These are the only two options that are worse after the BH
adjustment (Section~\ref{app:prep:overview}). Device-frame channels cost LR
\Rci{p1/d_lr_frame_device}~points, and restricting GBM to the gravity frame costs
\Rci{p1/d_gbm_frame_gravity}~points. All \R{p1/gbm_n_channels} channels in LR show no difference
(\Rci{p1/d_lr_frame_all}~points). The device-frame channels cannot act as a grip-to-hand shortcut here,
because the instructed grips are balanced across hands ($\chi^2$ test, $p$~=~\R{p1/grip_chi2_p}). They can,
however, fail on grips that differ from training: in the final-model selection the device-frame candidates
lose accuracy on the sideways grips and are removed by its grip test (Appendix~\ref{app:final}).

\paragraph{Window and time base.} The default window is the approach, time-normalised to
\R{p1/n_norm} samples. The options were the first half of the approach only, the approach plus the hold,
the approach with pre-onset context, and a fixed 3-s window. None changes accuracy beyond noise
(Table~\ref{tab:p1_decisions}). The largest gain forgone is \Rb{p1/d_mr_win_fixed_hi}~points (fixed
window, MR) and the largest loss compatible is \Rb{p1/d_lr_win_pre_lo}~points (pre-onset context, LR).
Decision: keep the approach. Production must segment the movement without a cue, and a window that needs
the hold or more context adds latency and depends more on the detector.

% ------------------------------------------------------------------------------------------------
\subsection{Robustness to segmentation errors}\label{app:prep:segmentation}

Production has no GO cue, so the detector of Appendix~\ref{app:prod} finds the approach with errors of tens
to hundreds of milliseconds. To set a requirement for that detector, models trained on the oracle bounds
(the approach bounds the app stored) were tested on bounds moved by independent uniform errors of up to
$\pm 50$, $\pm 100$ and $\pm 200$~ms, or with the onset 100~ms late (Figure~\ref{fig:P1_03_segmentation},
Table~\ref{tab:p1_robustness}).

No perturbation survives the BH adjustment in the robustness family. At $\pm 100$~ms the differences are
LR \Rci{p1/d_lr_bjit100} and GBM \Rci{p1/d_gbm_bjit100}~points. At $\pm 200$~ms GBM is nominally worse,
\Rci{p1/d_gbm_bjit200}~points (BH $q$~=~\R{p1/q_gbm_bjit200}). The largest loss compatible with the data
is therefore \Rb{p1/d_gbm_bjit200_lo}~points.

Training with jittered copies of the bounds (up to $\pm 100$~ms) was tried as a remedy. Against oracle-only
training at the same jittered test bounds it shows no difference: LR \Rci{p1/d_rem_lr_bjit100_trained}
and GBM \Rci{p1/d_rem_gbm_bjit100_trained}~points (BH $q$~=~\R{p1/q_rem_lr_bjit100_trained} and
\R{p1/q_rem_gbm_bjit100_trained} in the pairs family; Table~\ref{tab:app:prep-pairs}).
Decision: keep the oracle bounds in training. Treat bound errors as a requirement on the detector: at
errors up to $\pm 100$~ms no cost is detectable with \R{p1/n_persons} persons, and the data allow a loss of
at most \Rb{p1/d_lr_bjit100_lo}~points (LR). The errors here are synthetic and uniform; the real
detector's errors are evaluated in Appendix~\ref{app:prod}.

\begin{figure}[htbp]
\centering
\genfig{p1/P1_03_segmentation}
\caption{Robustness to segmentation errors: paired accuracy difference of a model trained on the oracle
bounds when its test bounds are moved by independent uniform errors of up to $\pm J$ (circles) or the
onset is 100~ms late (triangles). Diamonds: trained with jittered copies too, tested at $\pm 100$~ms and
judged against the oracle-trained model at the same bounds. 95\,\% CIs. Jitter training does not help; the
cost of bound errors is a requirement for the detector.}
\label{fig:P1_03_segmentation}\label{fig:app:prep-segmentation}
\end{figure}

\gentab{p1/p1_robustness}

% ------------------------------------------------------------------------------------------------
\subsection{Normalisation}\label{app:prep:norm}

The default feeds the channels in physical units. The options were a z-score per channel, a median and
interquartile-range scaling (both fitted on the mirror-augmented training rows, which keeps the mirror
symmetry) and a per-trial peak scaling (Figure~\ref{fig:P1_05_normalisation}). For LR the two per-channel
scalings are identical to the default, because LR standardises its own features (largest change of
\pright{}: \R{p1/maxdp_lr_norm_zscore}). No option differs from the default for any model. The widest CIs
belong to the per-trial peak scaling: LR \Rci{p1/d_lr_norm_trial_peak} and MR
\Rci{p1/d_mr_norm_trial_peak}~points. Decision: no normalisation. Physical units keep the speed and
amplitude of the movement; a per-trial peak scaling discards the amplitude.

\begin{figure}[htbp]
\centering
\genfig{p1/P1_05_normalisation}
\caption{Normalisation of the input channels against none (physical units), per panel model: paired
accuracy (left) and mean per-person log-loss (right) differences with 95\,\% CIs. Open squares: identical
predictions (LR standardises its own features). No option improves on physical units.}
\label{fig:P1_05_normalisation}\label{fig:app:prep-normalisation}
\end{figure}

% ------------------------------------------------------------------------------------------------
\subsection{Outliers, flags and exclusion rules}\label{app:prep:exclusions}

Two principles held throughout: evaluation trials are never dropped or relabelled, and any cleaning of the
training set must beat a size-matched control.

\paragraph{Training-side cleaning.} The options were winsorising at \R{p1/winsor_pct}\,\% and
$100 - \R{p1/winsor_pct}$\,\%, removing the trials of each quality flag (grip mismatch, false start,
segmenter disagreement above 100~ms, repeated samples, possible clipping) or of all flags from training,
and removing the swap-signature blocks (trials whose yaw sign suggests swapped hands). None helps
(Figure~\ref{fig:P1_01_forest}). Removing all flagged trials is nominally worse for GBM,
\Rci{p1/d_gbm_drop_all_flags}~points (BH $q$~=~\R{p1/q_gbm_drop_all_flags}). Against a control that removes
as many random trials per person and hand, it shows no difference: GBM
\Rci{p1/d_ctl_gbm_drop_all_flags} and LR \Rci{p1/d_ctl_lr_drop_all_flags}~points (pairs family,
Table~\ref{tab:app:prep-pairs}). The control alone costs LR \Rci{p1/d_lr_drop_all_flags_sizematch}~points
(BH $q$~=~\R{p1/q_lr_drop_all_flags_sizematch}). So the loss comes from the smaller training set, not from
the value of the flagged trials. Decision: keep flagged trials and swap blocks in training.

\paragraph{Scoring the removed trials.} An exclusion that removes scorable trials makes the headline
optimistic and the production gate too strict. Figure~\ref{fig:P1_09_exclusions} scores the removed and
flagged trials with the base models. The full table (\code{tables/p1/p1_subsets} of the analysis outputs)
is too long to reproduce; the main rows follow.
\begin{itemize}[nosep]
\item \emph{E4d} (an unbridged gap of more than \R{p1/max_gap_ms}~ms within \R{p1/gap_window_pad_s}~s of
  the movement) removes \R{p1/e4d_n} trials of \R{p1/e4d_n_persons} persons (\R{p1/e4d_n_android} Android,
  \R{p1/e4d_n_ios} iOS). The gap falls inside the approach in \R{p1/e4d_n_inapp} of them
  (\R{p1/e4d_n_persons_inapp} persons; median longest gap \R{p1/e4d_inapp_gap_median_ms}~ms) and only in
  the margins in \R{p1/e4d_n_margin}. With linear imputation, the in-approach trials score
  \Rci{p1/sub_lr_e4d_inapp_linear} (LR) and \Rci{p1/sub_gbm_e4d_inapp_linear} (GBM): compatible with the
  primary accuracy for LR and inconclusive for GBM. The margin-only trials score
  \Rci{p1/sub_lr_e4d_margin_linear} and \Rci{p1/sub_gbm_e4d_margin_linear}.
\item \emph{E4b/E4c} (shorter than \R{p1/min_duration_ms}~ms, or peak rotation below
  \R{p1/min_peak_dps}~deg/s) removes \R{p1/sub_gbm_e4bc_base_n} trials of \R{p1/sub_gbm_e4bc_base_np}
  persons. They score \Rci{p1/sub_lr_e4bc_base} (LR, inconclusive) and \Rci{p1/sub_gbm_e4bc_base} (GBM,
  below the primary accuracy).
\item \emph{Thresholds twice as strict} would remove \R{p1/sub_gbm_tight_fail_n} more trials of
  \R{p1/sub_gbm_tight_fail_np} persons, which score \Rci{p1/sub_gbm_tight_fail} (GBM): compatible with
  the primary accuracy, so stricter thresholds would discard scorable trials.
\item \emph{Subsets with fewer than 5 persons} get no CI; the values below are pooled over the trials of
  these few persons and are descriptive only. Grip-mismatch trials
  (\R{p1/sub_gbm_flag_grip_mismatch_n} trials, \R{p1/sub_gbm_flag_grip_mismatch_np} persons): LR
  \R{p1/sub_lr_flag_grip_mismatch}, GBM \R{p1/sub_gbm_flag_grip_mismatch}, MR
  \R{p1/sub_mr_flag_grip_mismatch}. Possible clipping (\R{p1/sub_gbm_flag_clip_candidate_n} trials,
  \R{p1/sub_gbm_flag_clip_candidate_np} persons): LR \R{p1/sub_lr_flag_clip_candidate}, GBM
  \R{p1/sub_gbm_flag_clip_candidate}, MR \R{p1/sub_mr_flag_clip_candidate}. Swap-signature blocks
  (\R{p1/sub_gbm_swap_n} trials of \R{p1/sub_gbm_swap_np} person): LR \R{p1/sub_lr_swap}, GBM
  \R{p1/sub_gbm_swap}, MR \R{p1/sub_mr_swap}. The swap blocks are a sensitivity analysis only; their labels
  were never changed.
\end{itemize}
Adding the excluded trials to training does not change the accuracy on the primary trials (E4d: LR
\Rci{p1/d_lr_e4d_trained}, GBM \Rci{p1/d_gbm_e4d_trained}~points; E4b/E4c: LR \Rci{p1/d_lr_excl_loose},
GBM \Rci{p1/d_gbm_excl_loose}~points).

\paragraph{Decision.} Keep E4b/E4c. They are justified a priori: a reach from arm's length lasts longer and
rotates the forearm, so these trials are fragments rather than reaches (Appendix~\ref{app:data}); the small
scored sample does not decide it. Keep E4d for comparability. There is no evidence that the imputed
in-approach trials score worse, but the sample is small, so production should flag an attempt with a gap
inside the movement rather than trust it.

\begin{figure}[htbp]
\centering
\genfig{p1/P1_09_exclusions}
\caption{Accuracy on the trials that an exclusion rule removes or a flag marks, scored by the base models
unless stated: E4d split by where the gap falls, with three imputations for gaps inside the approach;
E4b/E4c; trials that thresholds twice as strict would remove; flagged trials; the swap-signature blocks.
Dotted lines: all primary trials. 95\,\% CIs only for subsets of at least 5 persons (open markers: fewer, no
CI); n = trials, p.\ = persons. The removed trials are not clearly harder to score, except the E4b/E4c
fragments.}
\label{fig:P1_09_exclusions}\label{fig:app:prep-exclusions}
\end{figure}

% ------------------------------------------------------------------------------------------------
\subsection{Platform transfer}\label{app:prep:platform}

Android and iOS differ in sensors, rates, gaps and sensor fusion. A model that transfers only within a
platform would need training data from each. The options trained on the same platform only or on the other
platform only (Figure~\ref{fig:P1_08_platform}). Both also shrink the training set, so each was compared
with a control that removes as many training persons at random.

Against the pooled pipeline, GBM trained on the other platform is nominally worse,
\Rci{p1/d_gbm_platform_cross}~points (BH $q$~=~\R{p1/q_gbm_platform_cross}), and LR trained on the same
platform is nominally worse, \Rci{p1/d_lr_platform_same}~points (BH $q$~=~\R{p1/q_lr_platform_same}).
Against the size-matched controls, only GBM cross-platform training remains nominally worse,
\Rci{p1/d_ctl_gbm_platform_cross}~points (BH $q$~=~\R{p1/q_ctl_gbm_platform_cross} in the pairs family);
the other three pairs show no difference (Table~\ref{tab:app:prep-pairs}). Decision: pool both platforms.
Platform is a stratum, never a model input. The suggestive GBM loss is a reason to keep recruiting on both
platforms.

\begin{figure}[htbp]
\centering
\genfig{p1/P1_08_platform}
\caption{Platform transfer: accuracy on the persons of each platform (\R{describe/persons_android} Android,
\R{describe/persons_ios} iOS) when the model is trained on all other persons (pooled), on the same platform
only, or on the other platform only, each next to its size-matched control (open markers: as many training
persons drawn at random). LOSO; 95\,\% CIs. Only the difference from the control isolates the platform;
pooling is kept.}
\label{fig:P1_08_platform}\label{fig:app:prep-platform}
\end{figure}

% Compact version of the generated tables/p1/p1_pairs.tex, whose long row labels wrap into a table about a
% page tall: the same 16 comparisons (the arm accuracies and log-loss differences are in the generated file).
% To show the full table instead, replace this float by \gentab{p1/p1_pairs} and point the references to
% tab:p1_pairs.
\begin{table}[htbp]
\centering
\small
\caption{The pairs family: remedies against the problem they address, options against their size-matched
controls, and the controls against the current pipeline. Paired accuracy difference A minus B (points) with
its 95\,\% CI; $q$: BH within the \R{p1/fam_pairs_size} comparisons of the family. The arm accuracies and
log-loss differences are in \code{tables/p1/p1_pairs} of the analysis outputs.}
\label{tab:app:prep-pairs}
\begin{tabular}{@{}>{\raggedright\arraybackslash}p{6.4cm}lrrl@{}}
\toprule
A against B & Model & Difference [CI] & $q$ & Verdict \\
\midrule
\multicolumn{5}{@{}l}{\emph{Remedies}}\\
Jitter-trained vs oracle-trained, both tested at $\pm 100$~ms & LR & \Rcib{p1/d_rem_lr_bjit100_trained} &
  \R{p1/q_rem_lr_bjit100_trained} & \R{p1/v_rem_lr_bjit100_trained} \\
 & GBM & \Rcib{p1/d_rem_gbm_bjit100_trained} & \R{p1/q_rem_gbm_bjit100_trained} &
  \R{p1/v_rem_gbm_bjit100_trained} \\
E4d trials in training vs primary-only training, scored on primary and E4d trials & LR &
  \Rcib{p1/d_rem_lr_e4d_trained} & \R{p1/q_rem_lr_e4d_trained} & \R{p1/v_rem_lr_e4d_trained} \\
 & GBM & \Rcib{p1/d_rem_gbm_e4d_trained} & \R{p1/q_rem_gbm_e4d_trained} & \R{p1/v_rem_gbm_e4d_trained} \\
\midrule
\multicolumn{5}{@{}l}{\emph{Options against their size-matched controls}}\\
All flagged trials out of training vs as many random trials out & LR & \Rcib{p1/d_ctl_lr_drop_all_flags} &
  \R{p1/q_ctl_lr_drop_all_flags} & \R{p1/v_ctl_lr_drop_all_flags} \\
 & GBM & \Rcib{p1/d_ctl_gbm_drop_all_flags} & \R{p1/q_ctl_gbm_drop_all_flags} &
  \R{p1/v_ctl_gbm_drop_all_flags} \\
Same platform only vs as many random persons & LR & \Rcib{p1/d_ctl_lr_platform_same} &
  \R{p1/q_ctl_lr_platform_same} & \R{p1/v_ctl_lr_platform_same} \\
 & GBM & \Rcib{p1/d_ctl_gbm_platform_same} & \R{p1/q_ctl_gbm_platform_same} &
  \R{p1/v_ctl_gbm_platform_same} \\
Other platform only vs as many random persons & LR & \Rcib{p1/d_ctl_lr_platform_cross} &
  \R{p1/q_ctl_lr_platform_cross} & \R{p1/v_ctl_lr_platform_cross} \\
 & GBM & \Rcib{p1/d_ctl_gbm_platform_cross} & \R{p1/q_ctl_gbm_platform_cross} &
  \R{p1/v_ctl_gbm_platform_cross} \\
\midrule
\multicolumn{5}{@{}l}{\emph{Controls against the current pipeline}}\\
As many random trials out as the flagged ones & LR & \Rcib{p1/d_lr_drop_all_flags_sizematch} &
  \R{p1/q_lr_drop_all_flags_sizematch} & \R{p1/v_lr_drop_all_flags_sizematch} \\
 & GBM & \Rcib{p1/d_gbm_drop_all_flags_sizematch} & \R{p1/q_gbm_drop_all_flags_sizematch} &
  \R{p1/v_gbm_drop_all_flags_sizematch} \\
As many random persons as the same platform & LR & \Rcib{p1/d_lr_platform_same_sizematch} &
  \R{p1/q_lr_platform_same_sizematch} & \R{p1/v_lr_platform_same_sizematch} \\
 & GBM & \Rcib{p1/d_gbm_platform_same_sizematch} & \R{p1/q_gbm_platform_same_sizematch} &
  \R{p1/v_gbm_platform_same_sizematch} \\
As many random persons as the other platform & LR & \Rcib{p1/d_lr_platform_cross_sizematch} &
  \R{p1/q_lr_platform_cross_sizematch} & \R{p1/v_lr_platform_cross_sizematch} \\
 & GBM & \Rcib{p1/d_gbm_platform_cross_sizematch} & \R{p1/q_gbm_platform_cross_sizematch} &
  \R{p1/v_gbm_platform_cross_sizematch} \\
\bottomrule
\end{tabular}
\end{table}

% ------------------------------------------------------------------------------------------------
\subsection{Augmentation}\label{app:prep:augmentation}

Mirror augmentation adds the mirrored copy of every training trial with the other hand's label
(Equation~\eqref{eq:mirror}); it encodes the left-right symmetry of the task. The options were no mirror
augmentation, an added time warp (amplitude up to \R{p1/warp_max})~\cite{um2017augmentation}, and added
screen-away flips (GBM only). Without mirror augmentation, GBM accuracy shows no difference
(\Rci{p1/d_gbm_no_mirror}~points) but its log loss rises by \Rcib{p1/dll_gbm_no_mirror}. Screen-away
flips are nominally worse for GBM, \Rci{p1/d_gbm_screen_flip}~points (BH $q$~=~\R{p1/q_gbm_screen_flip}).
Time warping is nominally worse for LR, \Rci{p1/d_lr_time_warp}~points (BH $q$~=~\R{p1/q_lr_time_warp}),
and nominally lowers the GBM log loss. Decision: keep mirror augmentation alone.

% ------------------------------------------------------------------------------------------------
\subsection{Calibration of the panel's P(right)}\label{app:prep:calibration}

The product is a probability, so the panel's \pright{} was checked for calibration
(Figure~\ref{fig:P1_06_reliability}). Table~\ref{tab:p1_calibration} gives every calibrator,
Table~\ref{tab:p1_calibration_paired} its paired difference from the uncalibrated model, and
Table~\ref{tab:p1_calibration_platform} the result per platform.
Calibration was summarised by three numbers. The calibration slope is the coefficient of a logistic
regression of the hand on logit \pright{} (1 = calibrated, below 1 = over-confident). The intercept
(calibration-in-the-large) is its offset. The ECE was computed over \R{p1/ece_bins} equal-mass bins and
compared with the ECE of perfectly calibrated predictions on the same bins (\R{p1/ece_null_sims}
simulations). This null treats trials as independent, so it is supporting evidence only.

\paragraph{Remedies.} Four fixed calibrators (Platt~\cite{platt1999probabilistic},
temperature~\cite{guo2017calibration}, isotonic and beta) were fitted inside every outer training fold on
cross-fitted predictions of the training persons (\R{p1/inner_splits} grouped inner folds) and applied to
the held-out person. Comparing these \R{p1/n_fixed_calibrators} on the outer folds is exploratory. The
headline is the \emph{nested} choice: inside each training fold, grouped cross-validation chose among
\R{p1/n_auto_candidates} candidates (none, temperature, Platt, beta). These maps are strictly increasing, so
they keep the ranking within a fold. Isotonic is not a candidate because it ties predictions: it leaves
\R{p1/cal_gbm_isotonic_n_unique} distinct values of \R{p1/n_trials} for GBM and
\R{p1/cal_lr_isotonic_n_unique} for LR, and it costs LR ranking (AUC \Rcib{p1/calpd_lr_isotonic_roc_auc}).

\paragraph{GBM is over-confident and the nested calibrator corrects it.} Uncalibrated, its slope is
\Rcib{p1/gbm_slope} and its intercept \Rcib{p1/gbm_citl}. Its ECE is \R{p1/gbm_ece} against a null 95th
percentile of \R{p1/gbm_ece_null_q95} ($p$~=~\R{p1/gbm_ece_p}; BH $q$~=~\R{p1/caltest_gbm_ece_null_q} and
Holm $p$~=~\R{p1/caltest_gbm_ece_null_holm} in the calibration test family of \R{p1/caltest_family_size}
tests, Table~\ref{tab:p1_calibration_tests}). Using training persons only, the fitted temperature exceeds 1
in \R{p1/cal_gbm_t_above_1} of \R{p1/cal_gbm_t_n_folds} folds (median \R{p1/cal_gbm_t_median}, range
\R{p1/cal_gbm_t_min} to \R{p1/cal_gbm_t_max}). With the nested calibrator (temperature in
\R{p1/cal_gbm_auto_folds_temperature} folds, none in \R{p1/cal_gbm_auto_folds_none}, beta in
\R{p1/cal_gbm_auto_folds_beta}), the slope is \Rcib{p1/gbm_cal_slope}, the intercept
\Rcib{p1/gbm_cal_citl}, and the ECE \R{p1/gbm_cal_ece} against a null 95th percentile of
\R{p1/gbm_cal_ece_null_q95} ($p$~=~\R{p1/gbm_cal_ece_p}). The paired slope change is
\Rcib{p1/calpd_gbm_auto_slope} ($p$~\R{p1/caltest_gbm_slope_change_p}; BH $q$~=~\R{p1/caltest_gbm_slope_change_q}, Holm
$p$~=~\R{p1/caltest_gbm_slope_change_holm}). Its effect on ranking and accuracy is negligible: AUC
\Rcib{p1/calpd_gbm_auto_roc_auc}, accuracy \Rci{p1/calpd_gbm_auto_accuracy}~points. The log-loss gain,
\Rcib{p1/calpd_gbm_auto_log_loss}, is within noise at \R{p1/n_persons} persons. Decision: when a GBM gives the product
\pright{}, calibrate it with the nested calibrator.

\paragraph{The LR's calibrator hurts.} Uncalibrated, the LR slope is \Rcib{p1/lr_slope} (its CI covers 1)
and its intercept \Rcib{p1/lr_citl}; its ECE, \R{p1/lr_ece}, exceeds the null 95th percentile of
\R{p1/lr_ece_null_q95} ($p$~=~\R{p1/lr_ece_p}; BH $q$~=~\R{p1/caltest_lr_ece_null_q}, Holm
$p$~=~\R{p1/caltest_lr_ece_null_holm}). The nested calibrator (none in \R{p1/cal_lr_auto_folds_none} folds,
temperature in \R{p1/cal_lr_auto_folds_temperature}, Platt in \R{p1/cal_lr_auto_folds_platt}) gives a slope
of \Rcib{p1/lr_cal_slope} and an intercept of \Rcib{p1/lr_cal_citl}, but a larger ECE, \R{p1/lr_cal_ece},
against a null 95th percentile of \R{p1/lr_cal_ece_null_q95} ($p$~\R{p1/lr_cal_ece_p}). The slope moves by only
\Rcib{p1/calpd_lr_auto_slope} (BH $q$~=~\R{p1/caltest_lr_slope_change_q}, Holm
$p$~=~\R{p1/caltest_lr_slope_change_holm} in the calibration test family), while the paired log loss rises
by \Rcib{p1/calpd_lr_auto_log_loss} and the Brier score by \Rcib{p1/calpd_lr_auto_brier}. Decision: leave
the LR uncalibrated. Its deviations are not a monotone distortion that a rank-preserving map can remove.
The fixed isotonic map (exploratory) lowers the ECE to \R{p1/cal_lr_isotonic_ece} (null 95th percentile
\R{p1/cal_lr_isotonic_ece_null_q95}; slope \Rcib{p1/cal_lr_isotonic_slope}, intercept
\Rcib{p1/cal_lr_isotonic_citl}), but it ties predictions and costs ranking.

\paragraph{MR} carries its own inner Platt scaling. Its slope is \Rcib{p1/mr_slope}, its intercept
\Rcib{p1/mr_citl}, and its ECE \R{p1/mr_ece} against a null 95th percentile of \R{p1/mr_ece_null_q95}
($p$~=~\R{p1/mr_ece_p}, unadjusted: MR was not part of the calibration test family).

\paragraph{Per platform.} The pooled calibrator does not fully fix each platform
(Table~\ref{tab:p1_calibration_platform}). Uncalibrated GBM is most over-confident on iOS (slope
\Rcib{p1/cal_gbm_none_ios_slope}, intercept \Rcib{p1/cal_gbm_none_ios_citl}; ECE
\R{p1/cal_gbm_none_ios_ece}, null 95th percentile \R{p1/cal_gbm_none_ios_ece_null_q95}) and less so on
Android (slope \Rcib{p1/cal_gbm_none_android_slope}, intercept \Rcib{p1/cal_gbm_none_android_citl}; ECE
\R{p1/cal_gbm_none_android_ece}, null \R{p1/cal_gbm_none_android_ece_null_q95}). After the nested calibrator, the Android slope is
\Rcib{p1/cal_gbm_auto_android_slope} (intercept \Rcib{p1/cal_gbm_auto_android_citl}; ECE
\R{p1/cal_gbm_auto_android_ece}, null \R{p1/cal_gbm_auto_android_ece_null_q95}) and the iOS slope
\Rcib{p1/cal_gbm_auto_ios_slope} (intercept \Rcib{p1/cal_gbm_auto_ios_citl}; ECE
\R{p1/cal_gbm_auto_ios_ece}, null \R{p1/cal_gbm_auto_ios_ece_null_q95}). On each platform the ECE stays
slightly above its null band. With \R{describe/persons_ios} iOS persons this is a caveat, not a finding.

\begin{figure}[htbp]
\centering
\genfig{p1/P1_06_reliability}
\caption{Reliability of \pright{} for GBM (left) and LR (right), before calibration and after the
calibrator chosen inside every training fold (\R{p1/ece_bins} equal-mass bins; bars: 95\,\% CI of the
observed share of right hands; LOSO predictions). Legend: calibration slope with its CI, and the ECE with
the 95th percentile of the ECE of perfectly calibrated predictions on the same bins. The nested
calibrator brings the GBM slope to about 1; for LR it does not help.}
\label{fig:P1_06_reliability}\label{fig:app:prep-reliability}
\end{figure}

\gentab{p1/p1_calibration}
\gentab{p1/p1_calibration_paired}
\gentab{p1/p1_calibration_tests}
\gentab{p1/p1_calibration_platform}

% ------------------------------------------------------------------------------------------------
\subsection{Combining several attempts}\label{app:prep:aggregates}

The app may combine $k$ attempts of one person. Attempts of one person are correlated, so summing their
logits as if they were independent overstates the certainty, and averaging their probabilities understates
it. Four rules were compared on random $k$-subsets of one person and hand, for $k$ up to \R{p1/k_max}:
the mean of \pright{}, the mean of the logits, the sum of the logits, and the sum divided by a scale
\tauk{} fitted on the training persons of each fold, $\pright = \sigma(\sum_i z_i / \tauk)$, where $z_i$ is
the logit of attempt $i$ and $\sigma$ the logistic function (written $T_k$ in the figure and tables). The
same \R{p1/agg_gbm_n_persons} persons, those with at least \R{p1/k_max} trials of each hand, enter at every
$k$
(Figure~\ref{fig:P1_07_aggregates}, Tables~\ref{tab:app:prep-aggregates} and~\ref{tab:p1_aggregates_paired}).

For GBM at $k = 5$, the calibration slope is \Rcib{p1/agg_gbm_k5_sum_logit_slope} for the sum of the
logits (over-confident), \Rcib{p1/agg_gbm_k5_mean_prob_slope} for the mean of \pright{} and
\Rcib{p1/agg_gbm_k5_mean_logit_slope} for the mean of the logits (both under-confident). The scaled sum
gives \Rcib{p1/agg_gbm_k5_sum_logit_t_slope}, the closest to 1. Its median scale is
$\tau_5$~=~\R{p1/cal_gbm_t_sum_k5}, between 1 (the plain sum) and 5 (the mean of the logits). For LR the
scaled sum gives \Rcib{p1/agg_lr_k5_sum_logit_t_slope} ($\tau_5$~=~\R{p1/cal_lr_t_sum_k5}).
By log loss, the GBM rules cannot be told apart with \R{p1/agg_gbm_n_persons} persons: the scaled sum minus the sum is
\Rcib{p1/aggd_gbm_k5_sum_logit_ll}, minus the mean of the logits \Rcib{p1/aggd_gbm_k5_mean_logit_ll}, and
minus the mean of \pright{} \Rcib{p1/aggd_gbm_k5_mean_prob_ll}.
The aggregate analysis has no perfect-calibration null band for its ECE and no intercept, so this verdict
rests on the slope and the log loss only.

Decision: report the summed logits divided by \tauk{}, with \tauk{} refitted whenever the training data
change. It is the principled rule, not a demonstrated improvement over the mean of the logits.

\begin{figure}[htbp]
\centering
\genfig{p1/P1_07_aggregates}
\caption{Combining $k$ attempts of one person and hand (GBM; the same \R{p1/agg_gbm_n_persons} persons at
every $k$; every person and hand weighs the same). Left: ECE (no CI and no null band); right: log loss with
95\,\% CIs. Rules: mean of the probabilities, mean of the logits, sum of the logits, and that sum divided by
the training-fold scale \tauk{} ($T_k$). The scaled sum has the calibration slope closest to 1
(Table~\ref{tab:app:prep-aggregates}).}
\label{fig:P1_07_aggregates}\label{fig:app:prep-aggregates}
\end{figure}

% Compact version of the generated tables/p1/p1_aggregates.tex (40 rows; its ECE column has no
% perfect-calibration null band): calibration slopes and log losses only. To show the full table instead,
% replace this float by \gentab{p1/p1_aggregates} and point the references to tab:p1_aggregates.
\begin{table}[htbp]
\centering
\footnotesize\setlength{\tabcolsep}{4pt}
\caption{Calibration of $k$-attempt aggregates (\R{p1/agg_gbm_n_persons} persons, the same at every $k$;
every person and hand weighs the same): calibration slope at $k = 2$, 3 and 5 (CI only at $k = 5$) and log
loss at $k = 5$ with 95\,\% CIs. At $k = 1$ the rules coincide (slope \R{p1/agg_gbm_k1_mean_prob_slope}
for GBM, \R{p1/agg_lr_k1_mean_prob_slope} for LR). A slope above 1 is under-confident, below 1
over-confident. The full table, with accuracy, Brier score and ECE at every $k$, is
\code{tables/p1/p1_aggregates} of the analysis outputs.}
\label{tab:app:prep-aggregates}
\begin{tabular}{@{}llrrrr@{}}
\toprule
Model & Rule & Slope, $k = 2$ & Slope, $k = 3$ & Slope, $k = 5$ [CI] & Log loss, $k = 5$ [CI] \\
\midrule
GBM & mean of \pright{} & \R{p1/agg_gbm_k2_mean_prob_slope} & \R{p1/agg_gbm_k3_mean_prob_slope} &
  \Rcib{p1/agg_gbm_k5_mean_prob_slope} & \Rcib{p1/agg_gbm_k5_mean_prob_log_loss} \\
 & mean of logits & \R{p1/agg_gbm_k2_mean_logit_slope} & \R{p1/agg_gbm_k3_mean_logit_slope} &
  \Rcib{p1/agg_gbm_k5_mean_logit_slope} & \Rcib{p1/agg_gbm_k5_mean_logit_log_loss} \\
 & sum of logits & \R{p1/agg_gbm_k2_sum_logit_slope} & \R{p1/agg_gbm_k3_sum_logit_slope} &
  \Rcib{p1/agg_gbm_k5_sum_logit_slope} & \Rcib{p1/agg_gbm_k5_sum_logit_log_loss} \\
 & sum of logits / \tauk{} & \R{p1/agg_gbm_k2_sum_logit_t_slope} & \R{p1/agg_gbm_k3_sum_logit_t_slope} &
  \Rcib{p1/agg_gbm_k5_sum_logit_t_slope} & \Rcib{p1/agg_gbm_k5_sum_logit_t_log_loss} \\
\midrule
LR & mean of \pright{} & \R{p1/agg_lr_k2_mean_prob_slope} & \R{p1/agg_lr_k3_mean_prob_slope} &
  \Rcib{p1/agg_lr_k5_mean_prob_slope} & \Rcib{p1/agg_lr_k5_mean_prob_log_loss} \\
 & mean of logits & \R{p1/agg_lr_k2_mean_logit_slope} & \R{p1/agg_lr_k3_mean_logit_slope} &
  \Rcib{p1/agg_lr_k5_mean_logit_slope} & \Rcib{p1/agg_lr_k5_mean_logit_log_loss} \\
 & sum of logits & \R{p1/agg_lr_k2_sum_logit_slope} & \R{p1/agg_lr_k3_sum_logit_slope} &
  \Rcib{p1/agg_lr_k5_sum_logit_slope} & \Rcib{p1/agg_lr_k5_sum_logit_log_loss} \\
 & sum of logits / \tauk{} & \R{p1/agg_lr_k2_sum_logit_t_slope} & \R{p1/agg_lr_k3_sum_logit_t_slope} &
  \Rcib{p1/agg_lr_k5_sum_logit_t_slope} & \Rcib{p1/agg_lr_k5_sum_logit_t_log_loss} \\
\bottomrule
\end{tabular}
\end{table}

\gentab{p1/p1_aggregates_paired}

% ------------------------------------------------------------------------------------------------
\subsection{Accuracy by instructed condition}\label{app:prep:conditions}

Table~\ref{tab:p1_conditions} stratifies the GBM predictions by the instructed condition. Conditions were
never model inputs, no test was run, and the GBM designation is exploratory. Accuracy is
\Rci{p1/cond_speed_fast} at fast speed and \Rci{p1/cond_speed_slow} at slow speed, and
\Rci{p1/cond_tier_3} in protocol tier 3, the lowest of the four tiers. Strata with fewer than 5 persons have no CI. The
secondary analyses of speed and grip are in Appendix~\ref{app:sec}.

\gentab{p1/p1_conditions}

\FloatBarrier
% ------------------------------------------------------------------------------------------------
\subsection{Caveats}\label{app:prep:caveats}

\begin{itemize}[nosep]
\item The panel is small and fixed. Not every option ran on every model; MR ran only on the most
  informative options, to save compute.
\item With \R{p1/n_persons} persons, the resolution is a few points. ``No difference'' is absence of evidence; the lower
  CI bound is the largest loss the data allow.
\item Every choice was judged on the oracle bounds. The cue-free detector could change the ranking of the
  window options.
\item The segmentation errors are synthetic and uniform, not those of a real detector.
\item The ECE null treats trials as independent. The fixed calibrators are compared on the outer folds
  (exploratory, \R{p1/n_fixed_calibrators} configurations).
\item The GBM panel model uses device-frame channels. In the final-model selection such models lose
  accuracy on the sideways grips and are removed by the grip test (Appendix~\ref{app:final}), so the GBM
  verdicts transfer to the final model only where they do not depend on the device frame.
\item The variants ran the leak check on \R{p1/leak_sample_folds} seeded folds only; the base and
  nested-calibrated runs ran it on every fold.
\end{itemize}
